% Problem dossier: VI.05.P4
\subsection*{Problem VI.05.P4 --- Plane curves without a common component}
\label{prob:vi-05-p4}

\paragraph{Problem.}
\label{prob:vi05-p4}
Assume \(k\) is algebraically closed.  Let
\[
f,g\in k[x,y]
\]
have no common factor.

\begin{enumerate}
  \item Show that there exist
        \[
        u,v\in k[x,y]
        \]
        such that
        \[
        uf+vg
        \]
        is a nonzero polynomial in \(k[x]\).
  \item Assume \(f\) is irreducible.  Show that any proper algebraic subset of
        \[
        V(f)
        \]
        is finite.
  \item Deduce that two affine plane curves with no common irreducible component meet in
        only finitely many points.
  \item Prove that \(\mathbb A_k^2\) is quasi-compact in the Zariski topology.
\end{enumerate}

You may use the resultant in part (1), and Hilbert's basis theorem in part (4).

\ifdefined\FullProblemDossiers

\paragraph{Why this problem matters.}
This problem links elimination, irreducible curves, finite intersection theory, and
Noetherian/quasi-compact topology.  It is substantially more synthetic than the routine
exercises.

\paragraph{Useful ideas before starting.}
\begin{itemize}
\item Bézout in the PID \(k(x)[y]\), or the resultant with respect to \(y\);
\item a nonzero polynomial in one variable has finitely many roots;
\item irreducible plane curves correspond to prime principal ideals over an algebraically
closed field.
\end{itemize}

\paragraph{Guided hints.}
\textbf{Hint 1.} Regard \(f\) and \(g\) as polynomials in \(y\) over the field
\(k(x)\).  Their gcd is \(1\).

\textbf{Hint 2.} Clear denominators in a Bézout identity to produce a nonzero
\(h(x)\in(f,g)\).

\textbf{Hint 3.} If \(W\subsetneq V(f)\) is algebraic, choose a polynomial vanishing on
\(W\) but not on all of \(V(f)\), then apply part (1).

\textbf{Hint 4.} For quasi-compactness, reduce an arbitrary open cover to principal opens
and use finite generation of the unit ideal.

\paragraph{Solution.}
For part (1), regard \(f\) and \(g\) as polynomials in \(y\) with coefficients in \(k[x]\).
Because they have no common factor, their resultant with respect to \(y\),
\[
R(x)=\operatorname{Res}_y(f,g),
\]
is nonzero.  A standard Bezout identity for the resultant gives
\[
u(x,y)f(x,y)+v(x,y)g(x,y)=R(x)
\]
for suitable \(u,v\in k[x,y]\).

For part (2), let
\[
W\subsetneq V(f)
\]
be algebraic.  Since \(W\) is proper,
\[
I(W)\supsetneq I(V(f)).
\]
Assuming \(I(V(f))=(f)\), choose
\[
g\in I(W)\setminus(f).
\]
Because \(f\) is irreducible, \(f\) and \(g\) have no common factor.  By part (1),
\[
R(x)=uf+vg
\]
is a nonzero polynomial in \(k[x]\).

Every point of
\[
V(f,g)
\]
lies in \(V(R(x))\).  The polynomial \(R\) has finitely many roots
\[
a_1,\dots,a_N.
\]
Thus
\[
V(f,g)
\]
lies on finitely many vertical lines \(x=a_i\).

For a fixed \(a_i\), if both
\[
f(a_i,y)
\qquad\text{and}\qquad
g(a_i,y)
\]
were identically zero as polynomials in \(y\), then \(x-a_i\) would divide both \(f\) and
\(g\), contradicting the absence of a common factor.  Hence on each vertical line at least
one of the two restrictions is a nonzero polynomial in \(y\), so there are only finitely
many common zeros.  Therefore \(V(f,g)\), and hence \(W\), is finite.

For part (3), if two plane curves share no irreducible component, remove any repeated
factors and apply the same argument componentwise; their common zero set is finite.

For part (4), the polynomial ring \(k[x,y]\) is Noetherian, so the Zariski topology on
\(\mathbb A_k^2\) is Noetherian.  Every Noetherian topological space is quasi-compact.
Hence every Zariski open cover of \(\mathbb A_k^2\) has a finite subcover.

\paragraph{What happened mathematically.}
Elimination turns a two-variable intersection problem into a one-variable root problem.
This is the elementary precursor of resultants, elimination ideals, Bézout-type intersection
theory, and the Noetherian topology of affine algebraic sets.

\paragraph{Extensions.}
\begin{enumerate}
\item Compute an explicit resultant for two concrete plane curves.
\item Compare the affine statement with Bézout's theorem in projective space.
\item Find a pair of curves whose affine intersection is empty but whose projective closures
intersect at infinity.
\item Generalize the quasi-compactness argument to \(\mathbb A_k^n\).
\end{enumerate}

\paragraph{Provenance.}
Recovered from legacy `theory-of-differential-geometry-4.tex`, Problem 6.  The canonical version makes the algebraically closed-field hypothesis explicit.

\fi

